% DRAFT. Every number comes from bench/results/*.md in this repository;
% see paper/README.md for the mapping. Authors, affiliation and the
% benchmark's name are placeholders.
\documentclass[11pt]{article}
\usepackage[margin=1in]{geometry}
\usepackage[T1]{fontenc}
\usepackage{lmodern}
\usepackage{booktabs}
\usepackage{amsmath}
\usepackage{natbib}
\usepackage[hidelinks]{hyperref}
\usepackage{enumitem}
\setlist{nosep}

\title{Memory That Ages: A Longitudinal Benchmark for Agent Memory,\\
and What It Shows That Static Benchmarks Cannot}
\author{AUTHOR NAMES\\AFFILIATION\\\texttt{contact@example.org}}
\date{DRAFT --- November 2026}

\begin{document}
\maketitle

\begin{abstract}
Memory systems for language agents are compared on question-answering
benchmarks that load a fixed history and ask each question once. A
memory that lives with an agent for months does three things those
benchmarks cannot observe: it forgets, it learns from use which
memories matter, and it is asked about the same facts again. We
describe a small open-source memory system, Myelin, and a longitudinal
evaluation protocol that replays long conversations on a simulated
clock, asks each question some sessions after its evidence and
optionally again later, and scores retrieval without any language
model. On two standard benchmarks Myelin scores 90.1\% (LongMemEval-S,
\texttt{gpt-5-mini}, mean of four runs) and 92.6\% (LoCoMo), in the
range that published systems report; several mechanisms we added moved
those scores by less than their run-to-run spread. Under the
longitudinal protocol the same mechanisms separate clearly. Feedback
about which retrieved memories an answer used raises the share of
questions with all evidence in the top five results from 46.3\% to
52.7\%, and from 44.1\% to 53.7\% when questions are revisited.
Age-based forgetting alone hides a memory that a later question needs
for 8.6\% of questions (14.1\% of revisits); graded together with use
and with the use of related memories, that falls to 2.5\% (2.0\%).
We also report what did not work, the size of the noise in
model-judged scores, and the limits of our evidence. Code, data
loaders and every result file are public.
\end{abstract}

\section{Introduction}

An agent that works with a person for months needs a memory that
changes. New facts replace old ones. Most of what was said is never
needed again, and a little of it is needed repeatedly. Stores grow
until retrieving from them is slow or noisy unless something is
removed. Designs for agent memory have long reflected this: recency and
importance weighting in generative agents \citep{park2023generative},
an Ebbinghaus-style forgetting curve in MemoryBank
\citep{zhong2024memorybank,ebbinghaus1885}, tiered stores in MemGPT
\citep{packer2023memgpt}, temporal knowledge graphs
\citep{rasmussen2025zep} and linked notes \citep{xu2025amem}.

The benchmarks used to compare these systems do not change in the same
way. LongMemEval \citep{wu2025longmemeval} and LoCoMo
\citep{maharana2024locomo} each load a complete history into an empty
store and then ask a set of questions, each once. This is a good test
of whether the right text can be found and read. It says nothing about
a mechanism whose effect appears only as time passes or as questions
recur, because no time passes between writing and asking and no
question recurs.

We ran into this directly. While improving a memory system against
those two benchmarks we added forgetting that stops at quarantine,
feedback on which memories an answer used, and a health grade that
combines age, use and relationships between memories. None of these
could move a LongMemEval or LoCoMo score, by construction, and the
changes that could move it did so by amounts close to the difference
between two runs of an identical configuration. We therefore built a
protocol in which time passes, and report here what it shows.

Our contributions are:
\begin{itemize}
\item A longitudinal evaluation protocol for agent memory
  (Section~\ref{sec:protocol}): replay on a simulated clock, lagged and
  revisited questions, and three model-free measures (evidence in the
  top $k$, needed-but-hidden, and store still searchable).
\item A description of Myelin (Section~\ref{sec:system}), an MIT-licensed
  memory library on SQLite, with the measurements that led to each
  design choice, including the ones that were rejected.
\item Results under the protocol (Section~\ref{sec:results}) that
  separate mechanisms which static benchmarks cannot tell apart, and an
  account of how noisy model-judged scores are
  (Section~\ref{sec:static}).
\end{itemize}

\section{Related work}

\paragraph{Benchmarks.} LongMemEval \citep{wu2025longmemeval} asks 500
questions, each over its own history of roughly 115{,}000 tokens in the
S setting, with question types for single-session recall, multi-session
aggregation, temporal reasoning, knowledge updates and abstention.
LoCoMo \citep{maharana2024locomo} has ten conversations of 19 to 32
sessions spanning months, with about 150 questions each. Both are
static in the sense above. An independent audit of LoCoMo reports that
6.4\% of its answer key is wrong and that its commonly used judge
prompt accepts many wrong answers \citep{penfield2026locomo}.

\paragraph{Systems.} Published scores on these benchmarks are largely
self-reported, with different answering models, judges and prompts
\citep{chhikara2025mem0,rasmussen2025zep,hindsight2025,mastra2026om}.
We give them for orientation in Section~\ref{sec:static} and do not
treat them as a ranking.

\paragraph{Forgetting and activation.} Base-level activation in ACT-R
\citep{anderson2004actr} makes a memory's availability a function of
how often and how recently it was used, and spreading activation passes
some of it to associated memories. Our health grade
(Section~\ref{sec:health}) is in this tradition, in a much simpler
form, with one addition that matters for agents: a contradiction must
outweigh any amount of use.

\section{The system}
\label{sec:system}

Myelin stores one memory per message or fact in a single SQLite file.
A memory has text, a time from which it is valid, an optional time
from which it is not, metadata, and counters for how it has been used.

\paragraph{Retrieval.} Search combines BM25 \citep{robertson2009bm25}
over a stemmed full-text index with, optionally, cosine similarity over
embeddings from a small static model that runs on a CPU. The two
rankings are merged by reciprocal-rank fusion \citep{cormack2009rrf}.
A single fusion did not serve both the top of the list and its depth:
weighting the vector ranking double found the most evidence within a
6{,}000-token context on LongMemEval-S but was worse than lexical
search alone in the top ten on LoCoMo's short messages (39.6\% of
questions with all evidence in the top ten, against 54.9\%). Myelin
orders the first 20 results by an equal-weight fusion with a small rank
constant and the rest by the vector-heavy one, which matched the better
of the two at every depth measured (51.8\% in the top ten on LoCoMo,
77.5\% unchanged on LongMemEval-S, evidence within the budget
unchanged).

\paragraph{Context.} \texttt{recall} returns as many of the best
results as fit a token budget. A memory longer than a limit is shown
as an excerpt of the passages that share words with the query. On
LongMemEval-S, shortening the limit from 300 to 100 tokens raised the
share of questions with every evidence turn in a 6{,}000-token context
from 92.3\% to 97.0\%, because evidence is mostly in short user
messages and long assistant replies were using the budget.

\paragraph{Facts at write time.} Optionally, one model call per
conversation extracts short dated facts (``Ana adopted a dog on
2023-05-19'') that are indexed alongside the messages they came from;
the messages are still what is stored and shown. This raised LoCoMo
accuracy by about three points and made no measurable difference on
LongMemEval-S.

\paragraph{Feedback.} After an answer, \texttt{cite\_used} records
which of the shown memories it relied on. Those are counted as useful.
The others are counted as shown but not used, which lowers their rank
the next time they are relevant and leaves their health alone. The
answering model can supply the list: the context is numbered and the
prompt asks for a final line naming the entries used.

\subsection{One health grade}
\label{sec:health}

Forgetting is decided from a single number per memory,
$h = \sigma(r)$, with
\begin{align*}
r ={}& b - a(\text{age})
       - w_c\,(\text{contradictions})
       - w_q\,(1-\text{confidence})
       - w_i\,(1-\text{importance}) \\
     &+ \min\bigl(L,\;
        w_u \ln(1+u) + w_g \ln(1+s) + w_t\,\rho\bigr)
       - w_x \ln(1+x),
\end{align*}
where $a(\cdot)$ grows once a memory is more than thirty days old, $u$
is the number of times the memory was cited as used, $\rho$ decays with
the time since it was last used, $x$ counts cites in answers that were
wrong, and $s$ is support from related memories. Support flows from
each used memory to its neighbours: messages within three positions in
the same conversation, weighted by $1/(1+\text{distance})$, and its
eight nearest memories by embedding with cosine similarity at least
0.5, weighted by that similarity; in both cases scaled by
$\ln(1+u)$ of the source. A store with no cites has no support
anywhere, so the grade there depends on age, confidence and
importance alone.

The cap $L$ is set just below $w_c$. However much a memory, or its
neighbours, has been used, one contradiction outweighs it. This is
enforced at import time and by a test.

\paragraph{Forgetting is not deletion.} \texttt{prune} first marks
values that a later fact supersedes as no longer current, then
quarantines memories whose grade is below a threshold. Quarantined
memories leave search and stay in the store; they return if their
grade recovers, and are deleted only on explicit request.

\section{A longitudinal protocol}
\label{sec:protocol}

The protocol needs a dataset in which (i) content is dated over a long
period and (ii) each question is annotated with the items that answer
it. LoCoMo has both: sessions carry dates over 166 to 293 days, and
questions list their evidence utterances. We use its 1{,}531
non-adversarial questions that have evidence.

\paragraph{Replay.} For each conversation, a store is created with a
simulated clock. Sessions are written in order, each on its date.

\paragraph{Lagged questions.} Each question is scheduled a random
number of sessions (0 to $\ell$, default $\ell=6$) after the last
session containing its evidence, so that memories have aged by the
time they are needed. A question is asked by searching the store as it
then stands.

\paragraph{Revisits.} Optionally each question is asked again $n$ more
times, each a further 1 to $\ell$ sessions later. This models a user
returning to the same facts, which is the situation in which earlier
use of a memory can matter.

\paragraph{Arms.} The same schedule is run four ways: \emph{static}
(write and search only); \emph{decay} (\texttt{prune} after every
session); \emph{feedback} (after each question, \texttt{cite\_used}
with the question's evidence as the used set); and both.

\paragraph{Measures.} No language model is involved.
\begin{itemize}
\item \emph{Evidence in the top $k$}: the share of askings for which
  every evidence item is among the first $k$ results.
\item \emph{Needed but hidden}: the share of askings for which
  forgetting had already removed an evidence item from search.
\item \emph{Store still searchable}: the share of everything written
  that remains in search at the end.
\end{itemize}
The second and third are the two sides of forgetting. A system can
drive either to its best value trivially, by forgetting nothing or
everything; the protocol reports both, with ranking, so the exchange
is visible.

\paragraph{Feedback quality.} Using the evidence annotation as the
used set is an upper bound. In a separate held-out test on the same
data, letting \texttt{gpt-4o-mini} name the memories it used matched
that bound (47.8\% against 47.7\% of questions with all evidence in
the top five, against 42.4\% with no feedback). Feedback that gives
every shown memory one verdict for the whole answer gave no gain, and
neither did marking only the useful ones: the gain comes from moving
down memories that keep being retrieved and not used.

\section{What static benchmarks showed}
\label{sec:static}

Table~\ref{tab:static} gives accuracy on the two static benchmarks.
All Myelin runs use a 6{,}000-token context. LongMemEval-S is judged
by \texttt{gpt-4o} with the benchmark's prompts; LoCoMo by
\texttt{gpt-4o-mini} with the prompt that published scores use.

\begin{table}[t]
\centering
\small
\begin{tabular}{llr}
\toprule
Benchmark & Configuration & Accuracy \\
\midrule
LongMemEval-S & Release 1.4.0, \texttt{gpt-4o-mini} & 60.4 \\
 & Final, \texttt{gpt-4o} (one run) & 88.6 \\
 & Final, \texttt{gpt-5-mini} (four runs: 90.2, 91.0, 89.4, 89.8) & 90.1 \\
 & Evidence sessions given directly, \texttt{gpt-5-mini} & 93.4 \\
LoCoMo & Final, facts indexed, \texttt{gpt-5-mini} (one run) & 92.6 \\
 & Final, no facts, \texttt{gpt-5-mini} (one run) & 89.0 \\
\bottomrule
\end{tabular}
\caption{Myelin on static benchmarks (\%). ``Evidence sessions given
directly'' is a control that bypasses retrieval.}
\label{tab:static}
\end{table}

For orientation, published LongMemEval-S scores include 94.9 with
\texttt{gpt-5-mini} and 84.2 with \texttt{gpt-4o}
\citep{mastra2026om}, and 91.4 with a larger model
\citep{hindsight2025}; on LoCoMo, figures between 74.0 and 94.7 are
reported by the systems' authors, and an audit puts the ceiling of a
perfect system near 93.6 \citep{penfield2026locomo}. These use
different answering models, judges and prompts, and where third
parties have re-run a system the scores have differed widely. We make
no claim of rank.

Three observations from producing Table~\ref{tab:static} motivate the
rest of the paper.

\paragraph{The remaining gap is not in the memory.} With all evidence
in context for 96.4\% of LongMemEval-S questions, accuracy is 90.1\%;
with the evidence handed over directly it is 93.4\%. Most of what is
left is reading, concentrated in multi-session questions that require
counting items across sessions (about 80\% against 86.5\% in the
control). Five attempts to close it through retrieval (extra queries
written by a model, larger budgets, a model-based filter, write-time
facts and a stricter counting rule) did not.

\paragraph{The noise is as large as the effects.} Four runs of one
configuration ranged from 89.6\% to 93.0\%. Between two runs of an
identical setting about 40 of 500 answers change. One question type
with 30 questions moved between 76.7\% and 93.3\%. A change that moves
the score by a point cannot be told from no change without several
runs per setting, which few reports include. On LoCoMo, two runs of
one setting differed by about 50 answers each way out of 1{,}540.

\paragraph{Mechanisms over time are invisible.} Quarantine, feedback
and the health grade leave these scores untouched, since each question
meets a freshly written store. A change to ranking at the top of the
list that improved LoCoMo's top-ten evidence by twelve points left
accuracy within noise, because a 6{,}000-token context held the same
evidence either way.

\section{Results under the longitudinal protocol}
\label{sec:results}

All figures are means over three random question schedules on all ten
conversations, with hybrid search. The spread between schedules is
below half a point for ranking columns unless noted.

\begin{table}[t]
\centering
\small
\begin{tabular}{lrrrrr}
\toprule
 & \multicolumn{3}{c}{All evidence in} & Needed & Store still \\
 & top 5 & top 10 & top 20 & but hidden & searchable \\
\midrule
\multicolumn{6}{l}{\emph{Asked up to 6 sessions after the evidence}} \\
Static & 46.3 & 55.8 & 63.9 & 0.0 & 100.0 \\
Decay & 48.9 & 57.6 & 63.9 & 8.6 & 66.2 \\
Feedback & 52.7 & 60.6 & 64.0 & 0.0 & 100.0 \\
Decay + feedback & 52.8 & 60.8 & 64.9 & 2.5 & 86.2 \\
\midrule
\multicolumn{6}{l}{\emph{Asked up to 12 sessions after the evidence}} \\
Static & 44.8 & 54.7 & 62.8 & 0.0 & 100.0 \\
Decay & 46.9 & 55.5 & 61.4 & 11.8 & 66.2 \\
Feedback & 51.1 & 58.9 & 62.8 & 0.0 & 100.0 \\
Decay + feedback & 51.0 & 58.8 & 63.5 & 4.0 & 85.0 \\
\bottomrule
\end{tabular}
\caption{Longitudinal protocol, each question asked once (\%).}
\label{tab:once}
\end{table}

\begin{table}[t]
\centering
\small
\begin{tabular}{lrrrr}
\toprule
 & \multicolumn{2}{c}{All evidence in} & Needed & Store still \\
 & top 5 & top 10 & but hidden & searchable \\
\midrule
Static & 44.1 & 54.0 & 0.0 & 100.0 \\
Decay & 45.9 & 54.0 & 14.1 & 66.2 \\
Feedback & 53.7 & 59.5 & 0.0 & 100.0 \\
Decay + feedback & 53.8 & 59.6 & 2.0 & 89.6 \\
\bottomrule
\end{tabular}
\caption{Longitudinal protocol with two revisits: figures for repeat
askings only (\%).}
\label{tab:revisit}
\end{table}

\paragraph{Feedback.} Telling the store which memories were used is
the largest effect: 6.4 points of top-five evidence when each question
is asked once (Table~\ref{tab:once}) and 9.6 points on repeat askings
(Table~\ref{tab:revisit}). It carries over because questions about
the same events share evidence, and more so when the same question
returns.

\paragraph{Forgetting by age alone.} Decay removes a third of the
store from search and slightly sharpens the top of the list, since old
unused messages stop competing. It also hides a memory that a later
question needs for 8.6\% of questions, 11.8\% at the longer lag and
14.1\% of repeat askings. A sweep over decay speed and threshold
(36 settings on five conversations, checked on the other five) found
no setting that escaped this: without use signals the grade is a
function of age, the oldest memories go first whether or not they
matter, and roughly one asking in four hundred lost a needed memory
for each 1\% of the store hidden. Those two parameters are a dial, not
something to optimize.

\paragraph{Graded together.} With feedback, use keeps a memory
healthy, and support keeps its neighbours healthy. Needed-but-hidden
falls from 8.6\% to 2.5\% (11.8\% to 4.0\%; 14.1\% to 2.0\% on
revisits), with ranking equal to feedback alone or slightly better
and 10 to 15\% of the store still pruned.

\paragraph{What was tried and not kept.} A mild penalty on health for
being shown and not used pruned more (27\% of the store) and ranked a
point better in the top 20 at the shorter lag, and hid more needed
memories (6.2\% and 9.9\%); a stronger one collapsed the store. We
left health unaffected by unused showings. Counting related support in
ranking, as well as in health, made no difference at any weight: being
related to useful memories was a reason to keep a memory, not a reason
to rank it higher for a particular question. Adding the conversational
neighbours of top results to the context did not reliably help.

\section{Why a longitudinal benchmark is needed}

The argument does not depend on our system. Any memory that forgets,
consolidates, re-ranks from feedback or updates beliefs has behaviour
that is a function of history. A benchmark that writes once and asks
once evaluates that memory at a single point where none of the
behaviour has happened. Three consequences follow.

First, static scores cannot separate designs that differ only in how
they age. In Table~\ref{tab:once} the four arms differ by more than six
points in the top five; on LongMemEval and LoCoMo they are the same
system and receive the same score.

Second, static scores cannot show harm. Forgetting has a cost, and it
is the cost a deployed agent's user would notice: something they said
is gone when they ask for it. Needed-but-hidden measures that directly,
and it was large (8.6 to 14.1\%) for the simplest and most common
design, decay by age.

Third, static scores reward what is cheap to buy elsewhere. Our
accuracy on LongMemEval-S rose by about thirty points, of which the
largest steps came from the answering prompt and the answering model.
Those are real improvements to a product. They say little about
memory, and near the top of the range the differences between systems
are within the noise of the judge and the reader.

As agents are given longer-lived memory, the failures that matter
shift from ``could not find it'' to ``forgot it'', ``kept the wrong
version'' and ``keeps surfacing the same irrelevant thing''. Those are
properties of a trajectory. A benchmark should have one.

\section{Limitations}

\begin{itemize}
\item \textbf{One dataset.} The longitudinal results use LoCoMo only.
  Its conversations are social chat between two people; other domains
  may behave differently. Revisits repeat the same wording, which is
  easier than a paraphrase.
\item \textbf{Oracle feedback.} The feedback arms use evidence labels.
  A model's own account matched them in one separate test with one
  model; we have not run the full protocol with a model in the loop.
\item \textbf{Retrieval only.} The protocol scores whether evidence is
  retrievable, not whether answers are right. Our static results
  suggest answer accuracy is too noisy at this scale to resolve the
  effects, which is why we chose model-free measures, and it means the
  link to answer quality is argued and not measured.
\item \textbf{Small effects, few runs.} Several comparisons are within
  a point; schedules were repeated three times and model-judged runs
  between one and four times. We report spreads where we have them.
\item \textbf{Tuning.} The health weights were chosen on five
  conversations and checked on the other five of the same benchmark.
  That guards against fitting noise, not against fitting the dataset.
\item \textbf{Comparisons.} We did not run other systems. Published
  figures are quoted as their authors report them.
\item \textbf{Scale.} Stores here hold hundreds of memories. Separate
  timing on synthetic text shows retrieval latency growing linearly to
  about 45\,ms (lexical) and 83\,ms (hybrid) at 200{,}000 memories; the
  quality of forgetting at that scale is untested.
\end{itemize}

\section{Conclusion}

Benchmarks that ask each question once of a freshly written store
measure finding and reading. They are close to saturated by simple
retrieval and a strong reader, and noisy at the top. What distinguishes
memory systems that are meant to last (forgetting, promotion by use,
relationships between memories) appears only when time passes and
questions return. A protocol that replays dated conversations, lags and
repeats questions, and reports ranking together with what forgetting
cost, is cheap to run, needs no model, and separated mechanisms that
static benchmarks scored identically. We offer it, with a small system
built against it, as a starting point.

\paragraph{Availability.} The library, the harnesses for LongMemEval,
LoCoMo, cite feedback, scale and the longitudinal protocol, and every
result file cited here are in the project repository under the MIT
licence. LoCoMo and LongMemEval are distributed by their authors under
their own licences.

\bibliographystyle{plainnat}
\bibliography{refs}

\end{document}
